% ============================================================
% 第 8 章 SPH 流体与 GPU 粒子：Compute 管线三部曲
% 锚点：ParticleSystemGPU, FluidSystemGPU, SpatialHashGrid,
%       SPHKernelMath.h, GPUAsyncReadback + 14 个 sph/fluid/grid shader
% 时间线：29d41cc/beb5ac0(Phase 10) → 58a69bf(PCISPH)
%         → e46f088(SSFR) → 6cf61fb(Akinci 表面张力)
% ============================================================
\section{问题引入:一万颗粒子的 CPU 瓶颈}

粒子系统初版跑在 CPU 上，五万粒子时单帧模拟耗时超过 30ms——还没渲染就
已经跌破 30 FPS。而 GPU 正闲着。Phase 10 系列（\texttt{29d41cc} 基础设施
与 4-Pass 管线、\texttt{beb5ac0} 空间哈希网格与 SPH 流体）把整个模拟搬上
GPU Compute Shader，随后一路演进：PCISPH（\texttt{58a69bf}）、Screen-Space
Fluid Rendering（\texttt{e46f088}）、Akinci 表面张力（\texttt{6cf61fb}）。
本章是全书技术密度最高的一章，也是踩坑最多的一章。

\section{通用原理:SPH 与邻域搜索}

SPH（Smoothed Particle Hydrodynamics）把流体离散为携带质量/速度/密度的
粒子，任意场量用光滑核 $W(r,h)$ 加权求和近似：
$\rho_i = \sum_j m_j W(\lvert \mathbf{r}_i - \mathbf{r}_j\rvert, h)$，
压力由状态方程给出，加速度含压力梯度、粘性与外力。WCSPH 每步直接解压力；
PCISPH（Solenthaler 2009）则迭代预测位置修正密度误差，允许更大时间步。
本项目两者都实现，共享同一套核函数数学（SPHKernelMath.h，header-only 可测）。

朴素邻域查找是 $O(N^2)$——对每对粒子算距离。空间哈希把它降到近似
$O(N)$：粒子按格子哈希分桶，只查邻居桶。GPU 实现的经典三段式是
hash $\to$ prefix sum $\to$ scatter。

\section{本项目的设计与决策}

\subsection{Compute 4-Pass 粒子管线}

ParticleSystemGPU 把粒子生命周期拆成四个 compute dispatch：

\begin{itemize}
  \item \textbf{emit}：按发射参数生成新粒子写入粒子池；
  \item \textbf{simulate}：积分速度与位置（可挂接 SPH 力或 Bullet 刚体碰撞）；
  \item \texttt{render\_args}：在 GPU 上统计存活数并填充 Indirect Draw 参数块；
  \item \textbf{billboard}：顶点着色器按粒子属性展开面向相机的四边形。
\end{itemize}

关键决策是第 3 步：存活计数与绘制参数完全在 GPU 上生成，CPU 全程不知道
具体粒子数，DrawArraysIndirect 一条命令完成提交——CPU-GPU 数据往返被
压缩到零。VMware SVGA3D 不支持 indirect 时回退实例化绘制（renderer.md
规则记录的兼容性分支）。

\subsection{SpatialHashGrid:三段式与 binding 布局}

Engine/src/Renderer/SpatialHashGrid.h 是网格的完整自述：

\begin{lstlisting}
// Grid SSBOs —— binding 编号是跨 shader 的全局契约
Ref<ShaderStorageBuffer> m_CellHash;      // per-particle hash (binding 1)
Ref<ShaderStorageBuffer> m_CellCount;     // cell histogram   (binding 6)
Ref<ShaderStorageBuffer> m_CellStart;     // prefix sum 结果  (binding 5)
Ref<ShaderStorageBuffer> m_SortedIndices; // 重排索引         (binding 7)

void Build(uint32_t aliveCount, bool usePredictedPos = false);
// Vulkan 主帧 cmd 录制版本（避免 SingleTime stall，见 SPEC D-3）
void BuildVulkan(VkCommandBuffer cmd, uint32_t aliveCount,
                 bool usePredictedPos = false);
\end{lstlisting}

prefix sum 本身又是三个 pass（块内扫描 / 块和扫描 / 块偏移加回），pass 间
必须插 memory barrier——省略任何一个都会读到脏数据。binding 编号表是
14 个 sph/fluid/grid shader 共同遵守的全局布局契约，它的脆弱性本章踩坑录
有血泪案例。

\subsection{PCISPH 与表面张力:站在文献肩膀上}

PCISPH 的 $\delta$ 参数没有手抄论文数值，而是按 Solenthaler 2009 的公式
在初始化时自动推导（\texttt{00146b7}）；表面张力实现 Akinci 2013 的曲率
修正项，且 GLSL 与 CUDA 双实现同步落地（\texttt{6cf61fb}）。学术参数的
``自动推导优于魔法数字''在这里得到验证：改光滑核半径时一切自适应。

\subsection{SSFR:从粒子到水面}

粒子本身不像流体。SSFR（Screen-Space Fluid Rendering）管线用五个额外
pass 把粒子变成液体：thickness（厚度累积染色）$\to$ depth（椭圆 splat
写最小深度）$\to$ bilateral smooth（深度域平滑）$\to$ 从深度重建法线
$\to$ composite（折射+反射+泡沫合成进主场景）。它牺牲了物体空间一致性
换取巨大性能优势——屏幕空间的平滑天然只在可见处花钱。

\section{实现走读:异步回读与 fence}

CPU 需要知道存活数（如发射器停止条件），最初实现每帧
\texttt{glGetBufferSubData} 同步回读——GPU 必须等所有在途工作完成才能
服务这个读取，实测造成 28ms 的每帧 stall（修复提交 \texttt{01de4e3}）。
方案是 GPUAsyncReadback 抽象：N 个 host-visible 缓冲轮转，第 $k$ 帧发起
copy、读第 $k-N$ 帧的结果，fence 保证就绪。代价是数据延迟 N 帧——对
``上一帧存活数''这种用途完全无感。这次教训后来直接塑造了 Vulkan 路径的
D-11 决策（3 槽 ring + 独立 fence，见第 14 章）。

\section{踩坑记录}

这一章的坑值得单独成节——每个都是 GPU 并发世界的经典陷阱：

\begin{description}
  \item[SSBO slot 被 Grid.Build 覆盖] PCISPH 迭代中途调用 Grid.Build()
        重建网格，后者绑定的 SSBO 恰好占了 slot 1——预测位置缓冲被顶掉，
        流体整场不可见且不报任何错（\texttt{e63d819}）。排查花了三天，
        最终靠逐 pass 二分绑定才定位。\emph{OpenGL 全局 binding 空间里，
        ``谁绑了什么''是隐式契约}。
  \item[GL\_BLEND 泄漏] SSFR 深度 pass 继承了上游未关闭的混合状态，
        最小深度混合成了``平均深度''，水面法线重建全错（\texttt{b85f9d5}
        在深度与平滑通道显式禁 blend）。全局状态机的又一记闷棍。
  \item[CUDA/GLSL 一致性] 密度自贡献项遗漏、预测位置恢复时机、积分顺序
        三处差异让两后端结果分道扬镳（\texttt{dc72895} 三连修）。浮点
        求和顺序不同本属正常，但算法级差异必须消灭——这正是附录 A 容差
        协议存在的原因。
  \item[差一帧的 alive list] emit 之后不重建 alive list 就 simulate，
        新粒子晚一帧参与模拟（\texttt{547e3ad}）。GPU 数据结构没有
        ``自动维护的不变量''，每次写入都要想清楚下游谁在读。
  \item[cellSize=2h 的冗余] 早期邻居搜索半径设 2h，27 邻域扫描近半无效；
        收紧为 h 后消除冗余（\texttt{7c6c746}）。理论最优与实现常量
        要对齐验证。
\end{description}

\section{小结与延伸思考}

三部曲的主线：4-Pass 管线确立 CPU 极简参与；三段式网格把邻域查询摊平到
GPU；文献算法（PCISPH/Akinci）以``自动推导 + 双端一致''方式落地；SSFR
补上最后一公里的观感。性能上，这套管线支撑了附录 A 五万粒子的正式基准。

延伸思考：屏幕空间流体在粒子被遮挡或飞出屏幕时会丢失厚度信息（边缘
``融化''伪影）；物体空间方法（如 Marching Cubes 提取 mesh）能根治但代价
高昂。另一个值得探索的方向是把 PCISPH 迭代也做进 indirect 结构，进一步
减少 CPU 参与——目前迭代次数仍是启动时常量。
